% Parte cumulativa do relatório. Resultados oficiais inseridos após auditoria.
\section{Otimização Multiobjetivo}
\label{sec:multiobjetivo}

As buscas mono-objetivo fornecem referências para explorar compromissos entre
custo, desvio de $SiO_2$ e desvio de $Al_2O_3$. Nesta etapa são utilizados os
métodos soma ponderada e $\varepsilon$-restrito sobre os mesmos dados,
representação, construtiva e vizinhanças do problema original. Seja $X$ o
conjunto de composições que satisfazem integralidade, massas, compatibilidades,
disponibilidades e faixas químicas já formuladas. Ambos os métodos mantêm
todas essas restrições; as restrições agregadas de desvio não substituem as
faixas químicas individuais de cada pilha.

\subsection{Normalização com referências empíricas do TP1}

Seja $S_{\mathrm{ref}}$ o conjunto das 15 soluções finais viáveis das buscas
mono-objetivo corrigidas, com cinco sementes para cada objetivo. Os três
objetivos foram recalculados em cada uma dessas composições, incluindo os
valores cruzados dos objetivos que não eram o alvo da respectiva busca.
Para $j\in\{1,2,3\}$, adotam-se
\begin{align}
 a_j &= \min_{x\in S_{\mathrm{ref}}} f_j(x), &
 b_j &= \max_{x\in S_{\mathrm{ref}}} f_j(x), &
 g_j(x) &= \frac{f_j(x)-a_j}{b_j-a_j}.
\end{align}
Os valores $a_j$ e $b_j$ são extremos da amostra, não estimativas certificadas
do ideal ou do nadir global. Por isso, soluções posteriores podem apresentar
$g_j<0$ ou $g_j>1$; esses valores são mantidos sem truncamento. A implementação
rejeita referências não finitas ou com amplitude degenerada. Todas as
configurações e repetições utilizam a mesma referência congelada, com origem
das composições e identificador de integridade no arquivo de referência.

\begin{table}[H]
\centering
\caption{Extremos empíricos usados na normalização. Valores exibidos arredondados;
os cálculos utilizam a precisão completa da referência armazenada.}
\label{tab:referencias_tp2}
\begin{tabular}{lrr}
\toprule
Objetivo & $a_j$ & $b_j$ \\
\midrule
$f_1$ (R\$) & \num{60300000} & \num{65840000} \\
$f_2$ (p.p.$^2$) & \num{2.3727567955} & \num{6.7057453980} \\
$f_3$ (p.p.$^2$) & \num{1.4361688772} & \num{5.9951294968} \\
\bottomrule
\end{tabular}
\end{table}

\subsection{Formulação por soma ponderada}

Para cada vetor de pesos $w$, resolve-se heuristicamente
\begin{equation}
 (P_w)\qquad \min_{x\in X}\;\sum_{j=1}^{3}w_jg_j(x),
 \qquad w_j\geq 0,\quad \sum_{j=1}^{3}w_j=1.
\end{equation}
A malha contém os 15 vetores $(i/4,j/4,k/4)$ com $i,j,k$ inteiros não
negativos e $i+j+k=4$, acrescidos do centro $(1/3,1/3,1/3)$. Assim, cada
varredura resolve 16 configurações, incluindo os três vértices, preferências
intermediárias e pesos iguais. A resolução é uma escolha operacional para
explorar diferentes compromissos com esforço finito, e não uma exigência
numérica do enunciado. Pesos não normalizam os objetivos por si mesmos e não
substituem a verificação de viabilidade. Soma ponderada e GVNS não asseguram
a recuperação de toda a fronteira de um problema discreto.

\subsection{Formulação por $\varepsilon$-restrito}

O custo normalizado é mantido como objetivo principal:
\begin{equation}
 (P_{\varepsilon})\qquad
 \min_{x\in X}\;g_1(x)\quad
 \text{sujeito a}\quad g_2(x)\leq\varepsilon_2,\quad
 g_3(x)\leq\varepsilon_3.
\end{equation}
Cada limite pertence a $\{0,1/3,2/3,1\}$, de modo que são resolvidos todos
os 16 pares do produto cartesiano, não apenas os quatro pares da diagonal.
A malha varia independentemente os dois limites de qualidade agregada.
Um limite normalizado zero corresponde ao mínimo empírico do respectivo
desvio no TP1, e não a desvio bruto nulo. Tetos iguais a um também não
eliminam as restrições físicas, nem tornam o subproblema automaticamente viável.

\subsection{Comparação, inviabilidade e observação dos candidatos}

O núcleo comum utiliza avaliação lexicográfica para ambos os métodos TP2.
Candidatos estruturalmente inválidos são rejeitados antes do cálculo dos
objetivos. Entre candidatos estruturalmente válidos, a comparação prioriza
viabilidade, violação agregada, escalar e desempate determinístico.
A violação química soma os quadrados dos excessos além da tolerância,
divididos pela amplitude da respectiva faixa. No método $\varepsilon$-restrito,
somam-se também os excessos positivos dos dois limites normalizados.
Qualidade e restrições $\varepsilon$ usam tolerância absoluta $10^{-9}$;
a soma dos pesos é validada com tolerância absoluta $10^{-12}$.

Na soma ponderada, o desempate utiliza $(\sum_jg_j,g_1,g_2,g_3)$.
No $\varepsilon$-restrito, utiliza $(g_2+g_3,g_2,g_3)$, sem compensar
uma piora no custo principal. Empate completo mantém a incumbente.
Essa regra TP2 é distinta da penalidade histórica preservada para reprodução
do TP1: multiplicadores de penalidade não são referências de normalização.

A trajetória aceita, a melhor viável do subproblema e o arquivo não dominado
são objetos distintos. Cada candidato avaliado é observado antes do filtro
de aceitação, inclusive se for rejeitado ou se encerrar o orçamento. A melhor
viável é atualizada separadamente. A ausência de uma solução viável encontrada
é registrada explicitamente e não constitui prova de inviabilidade matemática.

\subsection{Integração do GVNS e critérios de parada}

\begin{algorithm}[H]
\caption{GVNS escalar com observação e orçamento comum}
\label{alg:gvns_multiobjetivo}
\begin{algorithmic}[1]
\Require Configuração escalar, referência congelada, semente e limites
\State Inicializar RNG próprio e contadores; melhor viável $\gets\varnothing$
\State Gerar, avaliar e observar a solução construtiva como incumbente
\While{há orçamento e os critérios externos permitem continuar}
  \State $k\gets1$; registrar se houve melhoria da incumbente
  \While{$k\leq K_{\max}$ e há orçamento}
    \State Perturbar a incumbente com SHAKE, observando cada avaliação
    \State Aplicar VND com N1/N2/N3 e comparar pela chave lexicográfica
    \If{o resultado melhora estritamente a incumbente}
      \State Atualizar incumbente; $k\gets1$
    \Else
      \State $k\gets k+1$
    \EndIf
  \EndWhile
  \State Atualizar histórico e contador de iterações sem melhoria
\EndWhile
\State \Return Melhor viável observada, ou ausência explícita, e registros
\end{algorithmic}
\end{algorithm}

O orçamento inclui a avaliação inicial e avaliações repetidas, inclusive
dentro de SHAKE e VND. Tentativas estruturalmente inválidas são contabilizadas
separadamente. O VND retorna à primeira vizinhança após melhoria enquanto
dispõe dos cinco reinícios; depois disso, permanece na vizinhança atual
enquanto houver melhoria e avança quando não há melhoria. O limite não
equivale a cinco melhorias totais.

Os parâmetros de referência são 200 iterações externas, 60 sem melhoria,
$K_{\max}=3$, 20 amostras e cinco reinícios VND. O piloto utiliza teto de
\num{10000} avaliações por subproblema para avaliar tempo e viabilidade.
Após o piloto, foi mantido o teto de \num{10000} avaliações por run,
registrado no protocolo antes do início da campanha oficial. As 32 buscas
ponderadas do piloto encontraram viável; no $\varepsilon$-restrito, foram
18 de 32, com as ausências concentradas nos mesmos sete pares mais apertados
nas duas sementes. Esses resultados e o tempo de execução fundamentaram
a decisão de orçamento; não demonstram inviabilidade matemática desses pares.
Os dados piloto permanecem separados dos resultados oficiais.
Os dois métodos recebem o mesmo teto;
o esforço efetivamente utilizado e o motivo de parada são registrados.

\subsection{Repetições, arquivo não dominado e reprodutibilidade}

Uma execução de uma configuração é denominada \emph{run}; uma varredura
percorre as 16 configurações de uma abordagem. O desenho oficial contém
cinco varreduras por método, identificadas pelas sementes-base
$42,123,456,789,1011$, totalizando $2\times5\times16=160$ runs.
Cada combinação de base, método e configuração recebe uma semente efetiva
derivada por SHA-256, evitando dependência da ordem de execução. O piloto
utiliza bases externas $90001$ e $90002$ e saída própria, sem misturar seus
dados com a coleta oficial.

O arquivo usa dominância estrita de Pareto nos três objetivos brutos:
$x$ domina $y$ quando $f_j(x)\leq f_j(y)$ para todo $j$, com desigualdade
estrita em pelo menos um objetivo. A filtragem ocorre em 3D antes das
projeções. Composições físicas iguais agregam todas as suas origens;
composições distintas com vetor igual podem permanecer. Arquivos separados
por repetição permitem examinar variabilidade, enquanto a união filtrada
por método reúne as alternativas não dominadas encontradas em suas varreduras.
Nenhum desses arquivos certifica a fronteira global.

As seis projeções requeridas são $f_1\times f_2$, $f_1\times f_3$ e
$f_2\times f_3$ para cada método. Cores e símbolos identificam as repetições,
com escalas comuns entre métodos para cada objetivo. As tabelas devem
contabilizar todas as runs, inclusive aquelas sem viável encontrada.

Manifestos registram dados, referência, configuração, fontes, versões e seeds.
A retomada revalida runs completas antes de reutilizá-las e recomeça uma run
interrompida; arquivos temporários não representam conclusão. O verificador
independente recalcula objetivos, restrições e proveniência das composições
e recompõe os arquivos não dominados. Alterações incompatíveis exigem uma
campanha em diretório novo. A análise comparativa final depende da auditoria
da coleta oficial; a integração reduzida demonstra funcionamento, não
superioridade experimental de uma abordagem.
